\documentclass[12pt,a4paper]{article}
\usepackage{geometry}
\geometry{left=2.5cm,right=2.5cm,top=2.0cm,bottom=2.5cm}
\usepackage[english]{babel}
\usepackage{amsmath,amsthm}
\usepackage{amsfonts}
\usepackage[longend,ruled,linesnumbered]{algorithm2e}
\usepackage{fancyhdr}
\usepackage{ctex}
\usepackage{array}
\usepackage{listings}
\usepackage{color}
\usepackage{graphicx}
\usepackage{float}

\begin{document}

\title{
  {\heiti《算法分析与设计》第 $5$ 次作业
    \footnote{要求：1、分析题请用书面化语言给出详细分析过程。}
    }
}
\date{}

\author{
姓名：\underline{你的名字}~~~~~~
学号：\underline{你的学号}~~~~~~
成绩：\underline{~~~~~~~~~~~~~~~~~~}
}

\maketitle

\noindent
\section*{\bf \color{red}{算法分析题}}
\vspace{10pt}

\noindent
{\bf 题目1：}用后进先出的$D$-搜索法找出$n$-皇后问题的一个解。(要求给出分析和伪代码)

\vspace{5pt}
\noindent
{\bf 答：}

\vspace{10pt}
\noindent
{\bf 题目2：}用回溯法设计一个给图 $G$($V$, $E$) 着色的算法。我们假定图是用邻接矩阵表示，而可用颜色的集合是 $C = \{1, 2, ..., m\}$，$m$ 可视为常数。我们要求把所有合法的着色全部输出。
\begin{enumerate}
	\item[(a)]  用子程序实现绑定函数 $Valid(k, c)$，用来检查树中 $(k-1)$ 层的一个点  $(x[1], x[2], ..., x[k-1])$ 的一个儿子。\\
	提示：假设图的顶点集合是 $V = \{v_1, v_2, ..., v_n\}$，邻接矩阵是$A$，$A[i, j] = 1$ 表示顶点 $v_i$ 和 $v_j$ 之间有边, $1 \leq i$ ,  $j \leq n$。我们用一个 $n$-元组 ($x[1], x[2], ..., x[n]$) 表示对这 $n$ 个顶点的着色，$x[i]$ 是对顶点 $v_i$ 的着色，$1 \leq i \leq n$。
	\item[(b)]  写出回溯算法的伪码。（需先写回溯法的递归形式函数供后面主程序调用）
\end{enumerate}

\vspace{5pt}
\noindent
{\bf 答：}

\vspace{10pt}
\noindent
{\bf 题目3：}用最大价值先出的分枝限界法找一个图的最大团。我们假定图$G$是用邻接矩阵表示的。(要求给出分析和伪代码)

\vspace{5pt}
\noindent
{\bf 答：}

\end{document}
